\section{The Rule Matcher}
\label{sec:rulematcher}

\subsection{Matching}

A rule set is an ordered list of compiled rules.  For a path, the matcher evaluates
each prefix of the path in turn (the folder itself, its parent, and so on down to the
path), and the set \emph{excludes} the path if any prefix is excluded: an excluded
folder excludes everything in it, and no later rule (not even a negation) can bring it
back, as in git.  Within one prefix the last matching rule wins, which is what makes
negation in the ignore file work.  The deny file has no negation, so it is a plain
union of rules.

Matching a path \emph{below} a folder ignores the prefixes at or above that folder.  It is
used for the hide rules while listing, so that entering a hidden folder on purpose shows
what is inside it.

\subsection{Compiling a rule}

A rule is translated to a regular expression, case-insensitive and with the dot matching
a newline (so that a name containing a newline cannot defeat \code{**}), ending in
\code{\textbackslash s*\$}, so that a name with trailing white space still matches.  The
translation follows the gitignore glob: \code{*} and \code{?} do not cross a
separator; \code{**} matches folders; a class \code{[...]} never matches a separator
(a negated class excludes it explicitly, and a \code{/} inside a class is an error); a
trailing \code{/} makes the rule apply to folders only.  A pattern is anchored to the home
folder if it begins with \code{\~{}/}, is absolute if it begins with \code{/}, is relative
to the home folder if it contains a \code{/} elsewhere, and matches at any depth
otherwise or after a leading \code{**/}.

\subsection{Name equality}
\label{sec:fold}

Both the pattern and the path are folded before they are compared (\code{Fold}): Unicode
normalization (NFC), then \emph{full} case folding, then NFC again.  The rules must
match names as the volume compares them, and on macOS the case-insensitive volume uses
full folding, under which the sharp s equals \code{ss}, the long s equals \code{s},
the Kelvin sign equals \code{k}, and the ligatures \code{ff}, \code{fi}, \code{fl},
\code{ffi}, \code{ffl} and \code{st} equal their spelled-out forms.  Simple case
folding, which was used at first, treated these as different names, so
\code{\~{}/.\ss h} opened \code{\~{}/.ssh} while the rule for \code{.ssh} did not match.  The
set of aliases was \emph{measured} on the volume, by creating every one- and
two-letter name and probing every character in the Unicode range, and is a test.
The client-side filter uses the equivalent lower-, upper-, lower-casing, and the tests
check that the two agree on the measured aliases.

Names are compared this way so that a deny rule over-matches, which is safe.  Where the
result is an \emph{allow} (whether a path is inside a root), identity is used instead
(Section~\ref{sec:roots}).

\subsection{Speed}

Listing a folder evaluates the rules once per entry.  A rule set of about thirty rules
took about 29~$\mu$s per entry, which would add three seconds to a folder of 100{,}000
entries, so two optimizations were made, both of which must give the same answer as the
plain loop:

\begin{itemize}
  \item A rule whose pattern contains no slash can match only the last component of a
        path, so it is tested against that component alone (short and anchored) instead of
        against the whole path at every level.
  \item For a set without negation, all rules of each kind are combined into one
        expression, asked first; only when it says a match is possible are the rules
        consulted one by one, to say \emph{which} rule matched (for \code{--debug}).
\end{itemize}

The cost is about 9~$\mu$s per entry, and a property test checks that the quick path
never changes an answer.

\subsection{Loading a rule file}

A rule file is written by a person in an editor, and a deny rule that loads but never
matches is the worst outcome.  The parser therefore either normalizes or refuses.  It
strips a leading byte order mark (which otherwise broke the first rule) and accepts any
line ending (a bare carriage return joined lines).  It refuses, naming the line: an
inline comment, \code{\~{}name}, a quoted pattern, an invisible character, a lookalike
of the slash, negation in the deny file, and the malformed patterns of the
functional specification (\fq{RUL-12}).  Whatever loads, works as written.

\subsection{The built-in rules}

The core deny rules are compiled into the binary, so they apply even without a deny
file.  The credential rules are not anchored to the current home folder, because the same
folders exist in backups, clones, other accounts' homes and under a wrong \code{HOME}.  If
a deny file exists it is authoritative, and \code{CoreMissing} compares the core rules with the
file's rules \emph{as names} to decide which are not in effect.  The list of credential
locations for the hard-link index (Section~\ref{sec:linkindex}) mirrors the credential entries of
the core rules, for the given home and every folder in \code{/Users} and \code{/home}
(\code{homeParents}).  It names literal folders, so for Flatpak, whose rules match any
application ID, it names the usual IDs.  The list holds the locations of macOS and of
Linux on both systems (Section~\ref{sec:decisions}).
